
The best-performing method for spurious correlation robustness achieves worst-group accuracy (WGA) of 0.91 on Waterbirds while leaving the ResNet-50 backbone numerically identical to the ERM baseline (layer4 cosine similarity = 1.000000 $\pm$ 1e-14 across all 3 seed pairs). This stands in structural contrast to GroupDRO, which achieves WGA = 0.88 by substantially modifying backbone weights (backbone/head L2 ratio = 6.47--7.03 across seeds). We study this divergence through a pre-registered diagnostic framework: background attribute linear probe accuracy on frozen layer4 features (sklearn L-BFGS, C=1e9, D=2048, full test set N=5,794) applied to 12 publicly available ResNet-50 checkpoints (3 seeds $\times$ 4 methods, izmailovpavel/spurious\_feature\_learning). We verify a three-step mechanistic chain for GroupDRO: minority group upweighting (5.01\% of training data) $\rightarrow$ backbone gradient propagation (backbone/head L2 ratio = 6.47--7.03) $\rightarrow$ significantly reduced background linear decodability in layer4 features (GroupDRO mean 0.9530 vs. ERM mean 0.9838; one-sided paired t-test p = 0.0039, Cohen's d = 6.48, n=3 seed pairs; CONFIRMED). Probe accuracy correlates negatively with WGA across methods (Pearson r = $-0.504$, n=9, SUGGESTIVE; ERM+GroupDRO subset r = $-0.755$, p = 0.041, CONFIRMED). Our results establish two mechanistically distinct WGA improvement pathways --- backbone-level spurious encoding reduction (GroupDRO) and head recalibration on unchanged backbone features (DFR) --- and provide a low-cost, reproducible diagnostic protocol for characterizing future robustification methods.

